% Figure: Neyman-Pearson rejection region under a Gaussian null and a
% Gaussian alternative, with Type I (alpha) and Type II (beta) areas
% labelled.
\begin{figure}[htbp]
\centering
\begin{tikzpicture}
\begin{axis}[
    width=12cm, height=6cm,
    xmin=-4, xmax=8,
    ymin=0, ymax=0.45,
    axis lines=left,
    xlabel={Test statistic $T$},
    ylabel={Density},
    samples=200,
    domain=-4:8,
    no markers,
    every axis plot/.append style={very thick},
    xtick={0, 2, 4}, xticklabels={$\theta_0$, $c$, $\theta_1$},
    ytick=\empty,
    legend style={font=\small, draw=none, fill=none, at={(0.98, 0.95)}, anchor=north east},
    clip=false,
]
\addplot[engblue]
  {1/sqrt(2*pi)*exp(-x*x/2)};
\addlegendentry{$H_0$: $T \sim \mathcal{N}(\theta_0, 1)$}

\addplot[engred, dashed]
  {1/sqrt(2*pi)*exp(-(x-4)*(x-4)/2)};
\addlegendentry{$H_1$: $T \sim \mathcal{N}(\theta_1, 1)$}

\addplot[engred!40, fill=engred!25, fill opacity=0.6,
         domain=-4:2, samples=80]
  {1/sqrt(2*pi)*exp(-(x-4)*(x-4)/2)} \closedcycle;

\addplot[engblue!60, fill=engblue!40, fill opacity=0.6,
         domain=2:8, samples=80]
  {1/sqrt(2*pi)*exp(-x*x/2)} \closedcycle;

\node[engblue!70!black] at (axis cs:3.0, 0.04) {\small $\alpha$};
\node[engred!70!black] at (axis cs:1.1, 0.04) {\small $\beta$};

\draw[engblack, thick] (axis cs:2, 0) -- (axis cs:2, 0.42);
\node[above, font=\small] at (axis cs:2, 0.42) {reject $H_0$};
\end{axis}
\end{tikzpicture}
\caption[Type I and Type II errors under simple hypotheses]{The
Neyman-Pearson framing of a hypothesis test. The threshold $c$
partitions the sample space into a rejection region (right of $c$)
and an acceptance region. The Type I error rate $\alpha$ is the
$H_0$-probability of falling in the rejection region; the Type II
error rate $\beta$ is the $H_1$-probability of falling in the
acceptance region. The two cannot be reduced simultaneously at fixed
sample size; the Neyman-Pearson lemma identifies the
likelihood-ratio test as the one with smallest $\beta$ at any given
$\alpha$ \cite[ch.~10]{text:wasserman2004}.}
\label{fig:vol02:ch14:test-regions}
\end{figure}
